% !TeX root = ../main.tex
% NOTE: ต้องมี \usepackage{float} ใน preamble ของ main.tex
%       เพื่อให้ [H] float specifier ทำงานได้

\section{การทดลองที่ 2: สร้าง Library HC-SR04 และวัด Measurement Noise Covariance}
% \textit{ครอบคลุม Criteria: Part 2 --- สร้าง Library HC-SR04, หาค่า $R$
% จาก Variance ของ Raw Data (ส่วนหนึ่งของ 5.5 คะแนน)}

%% ─────────────────────────────────────────────────────────────────────────────
\subsection{บทนำ}

ในระบบควบคุมอัตโนมัติและหุ่นยนต์ ข้อมูลที่ได้รับจากเซนเซอร์มักมีสัญญาณรบกวน
(Measurement Noise) แฝงอยู่เสมอ หากนำค่าดิบ (Raw Data) เหล่านี้ไปใช้งานโดยตรง
อาจทำให้ระบบตอบสนองผิดพลาดหรือขาดเสถียรภาพได้
ก่อนที่จะออกแบบ Kalman Filter ใดๆ จำเป็นต้องทราบค่า Measurement Noise
Covariance ($R$) ของเซนเซอร์ที่ใช้งานจริงก่อน เนื่องจาก $R$ เป็นพารามิเตอร์
ที่บอกให้ Filter ทราบว่าควรเชื่อถือการวัดจาก Sensor มากน้อยเพียงใด

การทดลองนี้จึงมีเป้าหมายในการ (1) พัฒนา Library สำหรับอ่านค่าระยะทาง
จาก HC-SR04 บน STM32 Nucleo G474RE และ (2) วัดค่า Measurement Noise
Covariance ของเซนเซอร์ที่ระยะต่างๆ อย่างเป็นระบบ โดยใช้วิธีการทางสถิติ
ที่เหมาะสม ผลลัพธ์ที่ได้จะนำไปใช้เป็นค่า $R$ สำหรับ Kalman Filter
ในการทดลองที่ 3 และ 5 ต่อไป

%% ─────────────────────────────────────────────────────────────────────────────
\subsection{วัตถุประสงค์}

\begin{itemize}
    \item เพื่อพัฒนา Library \texttt{hcsr04.c}/\texttt{hcsr04.h} สำหรับอ่านค่า
          Raw Data จากเซนเซอร์ HC-SR04 บน STM32 Nucleo G474RE
          ให้มีความถูกต้อง
    \item เพื่อสกัดค่า Measurement Noise Covariance ($R$) จากความแปรปรวน
          ของข้อมูลดิบที่ระยะต่างๆ โดยใช้ 5 Trials, 500 Samples ต่อ Trial
    \item เพื่อแสดงให้เห็นว่า Variance ของ HC-SR04 เพิ่มขึ้นตามระยะทาง
          และเลือกค่า $R$ ที่เหมาะสมสำหรับ Operating Range ของระบบ
\end{itemize}

%% ─────────────────────────────────────────────────────────────────────────────
\subsection{สมมติฐาน}

\begin{itemize}
    \item สัญญาณรบกวนของเซนเซอร์ HC-SR04 มีการกระจายตัวแบบ Gaussian
          รอบค่าจริง
    \item ระดับของ Variance จะเพิ่มขึ้นตามระยะห่างของเป้าหมาย
          เนื่องจาก Beam Divergence ของคลื่นอัลตราโซนิค
    \item การวัดซ้ำ 5 Trials ต่อระยะจะให้ค่าประมาณ $R$
          ที่เสถียร
\end{itemize}

%% ─────────────────────────────────────────────────────────────────────────────
\subsection{ตัวแปร}

\begin{itemize}
    \item \textbf{ตัวแปรอิสระ:} ระยะห่างของเป้าหมายที่ตั้งค่าไว้คงที่
          (10~cm, 20~cm, และ 30~cm)
    \item \textbf{ตัวแปรตาม:} ระยะทางดิบที่วัดได้ ($z_k$, หน่วย cm)
          และความแปรปรวนของสัญญาณ ($s^2$, หน่วย cm$^2$)
    \item \textbf{ตัวแปรควบคุม:} คาบเวลาการสุ่มตัวอย่าง $T_s = 60$~ms
          (\texttt{HAL\_Delay(60)} + ระยะเวลาอ่านเซนเซอร์),
          ทิศทางการจัดวางเป้าหมายให้ตั้งฉากกับเซนเซอร์,
          จำนวนตัวอย่าง (500 Samples ต่อ Trial, 5 Trials ต่อระยะ)
\end{itemize}

%% ─────────────────────────────────────────────────────────────────────────────
\subsection{เอกสารและงานวิจัยที่เกี่ยวข้อง}

\subsubsection{เซนเซอร์วัดระยะ HC-SR04}

\paragraph{แผนผังขาเชื่อมต่อ (Pin Diagram)}

โมดูล HC-SR04 มีขาเชื่อมต่อจำนวน 4 ขา ดังนี้

\begin{itemize}
    \item \textbf{VCC:} ขารับไฟเลี้ยง แนะนำให้ใช้ในช่วงแรงดัน 3.3--5~V~DC
    \item \textbf{TRIG:} ขารับสัญญาณควบคุม (Input) ต้องการพัลส์ TTL
          ที่มีความกว้างอย่างน้อย $10\,\mu$s เพื่อเริ่มต้นการส่งคลื่น
    \item \textbf{ECHO:} ขาส่งสัญญาณออก (Output) ความกว้างพัลส์ Echo
          ที่เป็น HIGH จะแปรผันตรงกับระยะทางที่วัดได้
          อยู่ในช่วง $100\,\mu$s ถึง 18~ms
    \item \textbf{GND:} ขาต่อกราวด์วงจร
\end{itemize}

\begin{figure}[htbp]
    \centering
    \includegraphics[width=0.70\textwidth]{exp2_hcsr04_pin.jpg}
    \caption{แผนผังขาเชื่อมต่อ (Pin Diagram) ของเซนเซอร์ HC-SR04}
    \label{fig:hcsr04_pin}
\end{figure}

\paragraph{Timing Diagram และลำดับการทำงาน}

หลักการทำงานของ HC-SR04 แบ่งออกเป็น 3 ช่วงหลัก

\begin{enumerate}
    \item \textbf{การสั่งงาน (Trigger Input):} ไมโครคอนโทรลเลอร์ส่งสัญญาณพัลส์
          HIGH ไปที่ขา TRIG เป็นเวลาอย่างน้อย $10\,\mu$s
    \item \textbf{การส่งคลื่น (Sonic Burst):} เซนเซอร์สร้างคลื่นอัลตราโซนิค
          ความถี่ 40~kHz จำนวน 8 ไซเคิล และยิงออกไปในอากาศ
    \item \textbf{การสะท้อนและอ่านค่า (Echo Pulse):} ขา ECHO
          เปลี่ยนสถานะจาก LOW เป็น HIGH ทันทีที่ส่งคลื่นเสร็จ
          และกลับมาเป็น LOW เมื่อรับสัญญาณสะท้อนกลับมาได้
\end{enumerate}

\textit{ข้อควรระวัง:} ต้องเว้นระยะห่างระหว่างรอบการวัดอย่างน้อย 60~ms
เพื่อป้องกันคลื่นสะท้อนตกค้างจากรอบก่อนหน้า

\begin{figure}[htbp]
    \centering
    \includegraphics[width=0.90\textwidth]{exp2_hcsr04_timing.png}
    \caption{Timing Diagram ของ HC-SR04 แสดงลำดับสัญญาณ Trigger,
             Sonic Burst, และ Echo Pulse}
    \label{fig:hcsr04_timing}
\end{figure}

\paragraph{การคำนวณระยะทางจากพัลส์ Echo}

ที่อุณหภูมิห้องประมาณ 20°C ความเร็วของเสียงในอากาศมีค่าเท่ากับ 343.4~m/s
หรือ 0.0343~cm/$\mu$s เนื่องจากเวลา $T$ ที่วัดได้คือเวลาเดินทางแบบไป-กลับ
ระยะทางจริงจึงต้องหารด้วย 2 ดังนี้

\begin{equation}
    s = \frac{v_{\mathrm{sound}} \times T}{2}
      = \frac{0.0343\,[\mathrm{cm/\mu s}] \times T\,[\mu\mathrm{s}]}{2}
      = \frac{T}{58.31}
      \approx \frac{T}{58} \; \text{(cm)}
    \label{eq:hcsr04_derived}
\end{equation}

การปัดค่า 58.31 เป็น 58 เป็นการประมาณเพื่อลดภาระการคำนวณใน
ไมโครคอนโทรลเลอร์

%% ──────────────────────────────────────────────────────────────────
\subsubsection{การ Configure STM32 Nucleo G474RE สำหรับ HC-SR04}

การอ่านค่าจาก HC-SR04 บน STM32 Nucleo G474RE ต้องการการตั้งค่าฮาร์ดแวร์
สองส่วนหลัก

\begin{itemize}
    \item \textbf{GPIO Configuration:} ขา TRIG ตั้งเป็น GPIO Output Push-Pull
          สำหรับส่งพัลส์ $10\,\mu$s และขา ECHO ตั้งเป็น GPIO Input
          สำหรับรับสัญญาณ
    \item \textbf{Timer Configuration (TIM1):} ใช้เป็น Free-Running Counter
          ที่ความละเอียด $1\,\mu$s/count โดยตั้ง Prescaler = 169
          จาก System Clock 170~MHz ทำให้ Timer นับที่ 1~MHz พอดี
          สามารถวัดความกว้างพัลส์ Echo ได้โดยตรงในหน่วย microseconds
\end{itemize}

\begin{table}[htbp]
    \centering
    \caption{การเชื่อมต่อ HC-SR04 กับ STM32 Nucleo G474RE}
    \label{tab:hcsr04_wiring}
    \begin{tabular}{llp{6cm}}
        \hline
        \textbf{HC-SR04 Pin} & \textbf{Nucleo G474RE Pin} & \textbf{หน้าที่} \\
        \hline
        VCC  & 5V        & ไฟเลี้ยง \\
        GND  & GND       & กราวด์ \\
        TRIG & PA9 (D8)  & GPIO Output --- ส่งพัลส์ $10\,\mu$s \\
        ECHO & PA8 (D7)  & GPIO Input --- รับพัลส์วัดระยะ \\
        \hline
    \end{tabular}
\end{table}

%% ──────────────────────────────────────────────────────────────────
\subsubsection{การคำนวณ Sample Variance เพื่อหาค่า $R$}

ค่า Measurement Noise Covariance $R$ ของ Kalman Filter คือ Variance
ของ Noise ของเซนเซอร์ ซึ่งประมาณได้จาก Sample Variance ของ Raw Data
ที่เก็บขณะวัตถุอยู่นิ่ง ดังสมการ

\begin{equation}
    s^2 = \frac{\displaystyle\sum_{i=1}^{N}(z_i - \bar{z})^2}{N - 1}
    \label{eq:exp2_sample_variance}
\end{equation}

โดยที่ $z_i$ คือค่าที่อ่านได้จากเซนเซอร์แต่ละครั้ง $\bar{z}$ คือค่าเฉลี่ย
และ $N$ คือจำนวนตัวอย่าง การหารด้วย $N-1$ แทน $N$ เรียกว่า
Bessel's Correction เพื่อให้ Sample Variance เป็น Unbiased Estimator
ของ Population Variance ที่แท้จริง

เมื่อมีการทำซ้ำหลาย Trials ค่า $R$ ที่ใช้จริงคือค่าเฉลี่ยของ
Sample Variance ข้าม $M$ Trials

\begin{equation}
    R = \bar{s}^2 = \frac{1}{M} \sum_{j=1}^{M} s_j^2
    \label{eq:exp2_mean_variance}
\end{equation}

โดยที่ $s_j^2$ คือ Sample Variance ของ Trial ที่ $j$
และ $M$ คือจำนวน Trials ทั้งหมด

\subsubsection{Kalman Filter ใน Continuous Domain}

ก่อนที่จะนำ Kalman Filter ไปใช้งานบนไมโครคอนโทรลเลอร์ในรูปแบบ Discrete
จำเป็นต้องเข้าใจรูปแบบ Continuous Domain ก่อน
Kalman--Bucy Filter คือรูปแบบ Continuous ของ Kalman Filter
ซึ่งประมาณค่าสถานะของระบบ LTI ที่อธิบายด้วยสมการ

\begin{align}
    \dot{x}(t) &= Ax(t) + Bu(t) + w(t) \label{eq:ct_state} \\
    z(t)        &= Hx(t) + v(t)          \label{eq:ct_obs}
\end{align}

โดยที่ $w(t) \sim \mathcal{N}(0, Q_c)$ คือ Process Noise
และ $v(t) \sim \mathcal{N}(0, R)$ คือ Measurement Noise
สมการ Kalman--Bucy Filter ใน Continuous Domain มีรูปแบบ

\begin{equation}
    \dot{\hat{x}}(t) = A\hat{x}(t) + Bu(t) + K_f(t)\bigl(z(t) - H\hat{x}(t)\bigr)
    \label{eq:kalman_bucy}
\end{equation}

โดยที่ $K_f(t) = P(t)H^TR^{-1}$ คือ Kalman Gain แบบ Continuous
และ $P(t)$ คือ Error Covariance 

\begin{equation}
    \dot{P}(t) = AP(t) + P(t)A^T + Q_c - P(t)H^TR^{-1}HP(t)
    \label{eq:riccati}
\end{equation}

อย่างไรก็ตาม การ Implement สมการเหล่านี้บนไมโครคอนโทรลเลอร์โดยตรงไม่สามารถ
ทำได้ เนื่องจากระบบฝังตัวทำงานในเวลาแบบ Discrete จึงต้องทำการ Discretize
สมการก่อนนำไปใช้งาน ดังที่จะอธิบายในหัวข้อถัดไป

\subsubsection{การ Discretize สมการ Kalman Filter}

การแปลงสมการ Kalman Filter จาก Continuous Domain ไปเป็น Discrete Domain
ทำได้โดยใช้วิธี Zero-Order Hold (ZOH) หรือ Euler Forward
สำหรับ Scalar KF ที่ใช้กับ HC-SR04 ในการทดลองนี้
กำหนดพารามิเตอร์ของระบบดังนี้

\begin{itemize}
    \item ระบบมีสถานะเดียวคือระยะทาง $x_k$ (Scalar)
    \item ไม่มี Control Input ดังนั้น $B = 0$
    \item Sensor วัดสถานะได้โดยตรง ดังนั้น $H = 1$
    \item ระบบสมมติให้สถานะไม่เปลี่ยนแปลงระหว่าง Step
    ดังนั้น $A_c = 0$ ซึ่งนำไปสู่ $F = e^{A_c T_s} = 1$
\end{itemize}

เมื่อใช้ Euler Forward Approximation $F \approx I + A_c T_s$
และ $A_c = 0$ สำหรับระบบ Static Model จะได้

กระบวนการ Discretize เริ่มจาก Continuous-Domain State Equation:

\begin{equation}
    \dot{x}(t) = A_c\,x(t) + B_c\,u(t) + w(t)
    = 0 \cdot x(t) + 0 \cdot u(t) + w(t)
    \label{eq:scalar_continuous}
\end{equation}

ซึ่งสะท้อนว่าในแบบจำลอง Static Model สถานะจริงไม่มี
แนวโน้มเปลี่ยนแปลงเอง ความเปลี่ยนแปลงที่เกิดขึ้นถูกจำลอง
ผ่าน Process Noise $w(t)$ เท่านั้น

จากนั้นใช้ Zero-Order Hold (ZOH) เพื่อแปลง $A_c$ ไปเป็น
Discrete State Transition Matrix $F$:

\begin{equation}
    F = e^{A_c T_s}
      = e^{0 \cdot T_s}
      = e^{0}
      = 1
    \label{eq:scalar_F_derive}
\end{equation}

และ Discrete Input Matrix:

\begin{equation}
    B_d = \int_0^{T_s} e^{A_c \tau}\,B_c\,d\tau
        = \int_0^{T_s} 1 \cdot 0\,d\tau
        = 0
    \label{eq:scalar_Bd_derive}
\end{equation}

ผลลัพธ์ที่ได้คือ Discrete State-Space Model สำหรับ Scalar KF:

\begin{align}
    x_k      &= F\,x_{k-1} + B_d\,u_{k-1} + w_{k-1}
               = x_{k-1} + w_{k-1}
               \label{eq:scalar_state_disc} \\
    z_k      &= H\,x_k + v_k
               = x_k + v_k
               \label{eq:scalar_obs_disc}
\end{align}

ซึ่งสอดคล้องกับสมมติฐานที่ว่าระยะทางที่แท้จริงไม่เปลี่ยนแปลง
ระหว่าง Sampling Period และ Sensor วัดสถานะได้โดยตรงแต่มี Noise

%% ─────────────────────────────────────────
\subsubsection{สรุปสมการคณิตศาสตร์ของ Discrete Kalman Filter}

ก่อนที่จะนำอัลกอริทึมไปเขียนเป็นโปรแกรมลงบนไมโครคอนโทรลเลอร์ STM32 ขอสรุปสมการหลักทั้ง 5 ขั้นตอนของ Discrete Kalman Filter ซึ่งแบ่งการทำงานออกเป็น 2 ระยะ (Phase) อย่างเป็นวงรอบ (Recursive) ได้แก่ ระยะพยากรณ์ (Prediction Step) และระยะปรับปรุง (Update Step) ดังนี้



\textbf{1. ระยะพยากรณ์ (Prediction Step / Time Update)}
ขั้นตอนนี้ทำหน้าที่คาดการณ์สถานะล่วงหน้าหนึ่งสเต็ปเวลา โดยอาศัยแบบจำลองทางคณิตศาสตร์ (State-Space Model) ของระบบ
\begin{align}
    \hat{x}_{k|k-1} &= A \hat{x}_{k-1|k-1} + B u_k \label{eq:kf_predict_state} \\
    P_{k|k-1} &= A P_{k-1|k-1} A^T + Q \label{eq:kf_predict_cov}
\end{align}

\textbf{2. ระยะปรับปรุง (Update Step / Measurement Update)}
ขั้นตอนนี้ทำหน้าที่นำข้อมูลจริงที่วัดได้จากเซนเซอร์ ณ เวลานั้น มาปรับแก้ค่าที่พยากรณ์ไว้ให้มีความแม่นยำยิ่งขึ้น โดยให้น้ำหนักความเชื่อถือผ่านค่า Kalman Gain
\begin{align}
    K_k &= P_{k|k-1} H^T (H P_{k|k-1} H^T + R)^{-1} \label{eq:kf_kalman_gain} \\
    \hat{x}_{k|k} &= \hat{x}_{k|k-1} + K_k (z_k - H \hat{x}_{k|k-1}) \label{eq:kf_update_state} \\
    P_{k|k} &= (I - K_k H) P_{k|k-1} \label{eq:kf_update_cov}
\end{align}

รายละเอียดของตัวแปรในสมการมีดังนี้:
\begin{itemize}
    \item $\hat{x}$: สถานะของระบบ (State Estimate) เช่น ตำแหน่ง หรือ ความเร็ว
    \item $P$: ความแปรปรวนร่วมของความคลาดเคลื่อน (Error Covariance) สะท้อนถึงความไม่แน่นอนของการประมาณค่า
    \item $A, B, H$: เมทริกซ์พลวัตของระบบ (System Dynamics), เมทริกซ์อินพุต (Input Matrix), และเมทริกซ์สังเกตการณ์ (Measurement Matrix)
    \item $Q, R$: ค่าความแปรปรวนของสัญญาณรบกวนในกระบวนการ (Process Noise Covariance) และจากการวัด (Measurement Noise Covariance)
    \item $K_k$: อัตราขยายคาลมาน (Kalman Gain) ซึ่งเป็นตัวกำหนดน้ำหนักระหว่างแบบจำลองกับเซนเซอร์
    \item $z_k$: ค่าที่อ่านได้จริงจากเซนเซอร์ (Measurement)
\end{itemize}

\textit{หมายเหตุ:} สำหรับระบบแบบ Scalar (1 มิติ) เมทริกซ์ทั้งหมดจะลดรูปลงเป็นค่าคงที่ (Scalar value) การดำเนินการทรานสโพส ($^T$) จะไม่มีผล และการอินเวอร์สเมทริกซ์ ($^{-1}$) จะถูกแทนที่ด้วยการหารทางคณิตศาสตร์ทั่วไป

โค้ดด้านล่างแสดงการ Implement Scalar KF ในภาษา C บนไมโครคอนโทรลเลอร์ STM32 Nucleo G474RE โดยแต่ละขั้นตอนสอดคล้องกับสมการ~\eqref{eq:kf_predict_state}--\eqref{eq:kf_update_cov} อย่างสมบูรณ์

\subsubsection{การ Implement Scalar Kalman Filter บน STM32}

โค้ดด้านล่างแสดงการ Implement Scalar KF ใน C บน STM32 Nucleo G474RE
โดยแต่ละขั้นตอนสอดคล้องกับสมการ~\eqref{eq:kf_predict_state}--\eqref{eq:kf_update_cov}

\begin{lstlisting}[language=C, caption={Scalar Kalman Filter — kalman.c},
                   label={lst:scalar_kf}]
typedef struct {
    float x_est;   // Estimated state
    float P;       // Error covariance
    float Q;       // Process noise covariance
    float R;       // Measurement noise covariance
} KalmanScalar_t;

void Kalman_Init(KalmanScalar_t *kf, float Q, float R, float x0)
{
    kf->x_est = x0;
    kf->P     = 1.0f;
    kf->Q     = Q;
    kf->R     = R;
}

float Kalman_Update(KalmanScalar_t *kf, float z)
{
    // Prediction
    float x_prior = kf->x_est;
    float P_prior = kf->P + kf->Q;

    // Update
    float K = P_prior / (P_prior + kf->R);
    kf->x_est = x_prior + K * (z - x_prior);
    kf->P = (1.0f - K) * P_prior;

    return kf->x_est;
}
\end{lstlisting}

\noindent ใน Main Loop การใช้งาน Filter มีโครงสร้างดังนี้

\begin{lstlisting}[language=C, caption={main.c — Main Loop สำหรับ Scalar KF},
                   label={lst:main_kf}]
KalmanScalar_t kf;
float raw_cm, filtered_cm;

HAL_TIM_Base_Start(&htim1);
HCSR04_Init(&htim1);
raw_cm = HCSR04_Read();

Kalman_Init(&kf, Q_VALUE, R_VALUE, raw_cm);

while (1) {
    raw_cm = HCSR04_Read();

    if (raw_cm > 0.0f) {
        filtered_cm = Kalman_Update(&kf, raw_cm);
        printf("%lu,%.4f,%.4f\r\n", HAL_GetTick(), raw_cm, filtered_cm);
        fflush(stdout);
    }
    HAL_Delay(60);
}
\end{lstlisting}

%% ─────────────────────────────────────────────────────────────────────────────
\subsection{ขั้นตอนการทดลอง}

\begin{enumerate}
    \item พัฒนา Library \texttt{hcsr04.c} สำหรับเชื่อมต่อ HC-SR04
          กับ Nucleo G474RE โดยใช้ Timer TIM1 ที่ความละเอียด $1\,\mu$s
          และประมวลผลระยะทางเป็นชนิดข้อมูล \texttt{float}
    \item ต่อสาย HC-SR04 กับ Nucleo G474RE ตามตารางที่~\ref{tab:hcsr04_wiring}
    \item ส่งข้อมูลระยะทางดิบผ่าน LPUART1 (Baud rate 115200)
          ออกมาที่คอมพิวเตอร์เป็น CSV format
    \item วางวัตถุแบนตั้งฉากกับเซนเซอร์ที่ระยะเป้าหมาย 10~cm แล้วบันทึก
          500 ตัวอย่าง ทำซ้ำ 5 Trials โดยในแต่ละ Trial ปล่อยให้ระบบ
          นิ่งก่อนเริ่มบันทึกข้อมูล
    \item ทำซ้ำข้อ 4 ที่ระยะ 20~cm และ 30~cm
    \item นำข้อมูลเข้า MATLAB คำนวณ Sample Variance ของแต่ละ Trial
          ตามสมการ~\eqref{eq:exp2_sample_variance} และหาค่าเฉลี่ย
          ตามสมการ~\eqref{eq:exp2_mean_variance}
    \item เลือกค่า $R$ ที่สอดคล้องกับ Operating Range ของระบบ
          Mass-Spring-Damper
\end{enumerate}

%% ─────────────────────────────────────────────────────────────────────────────
\subsection{การพัฒนา Library HC-SR04}

\subsubsection{โครงสร้าง Library (\texttt{hcsr04.h} / \texttt{hcsr04.c})}

Library ประกอบด้วยสองฟังก์ชันหลัก ได้แก่
\texttt{HCSR04\_Init()} สำหรับตั้งค่า Timer และ GPIO
และ \texttt{HCSR04\_Read()} สำหรับวัดระยะจริง
โดยกำหนดค่าคงที่ไว้ใน \texttt{hcsr04.h}

\begin{lstlisting}[language=C,
    caption={\texttt{hcsr04.h} --- ค่าคงที่และ Function Prototype},
    label={lst:hcsr04_h}]
#define HCSR04_TRIG_PIN              GPIO_PIN_9
#define HCSR04_TRIG_PORT             GPIOA

#define HCSR04_ECHO_PIN              GPIO_PIN_8
#define HCSR04_ECHO_PORT             GPIOA

#define HCSR04_ECHO_HIGH_TIMEOUT_MS  30U
#define HCSR04_ECHO_LOW_TIMEOUT_MS   50U
#define HCSR04_TIMEOUT_VALUE         -1.0f
#define HCSR04_MIN_INTERVAL_MS       60U

void  HCSR04_Init(TIM_HandleTypeDef *htim);
float HCSR04_Read(void);
\end{lstlisting}

\subsubsection{การทำงานของ \texttt{HCSR04\_Read()}}

ฟังก์ชัน \texttt{HCSR04\_Read()} คืนค่าระยะทางในหน่วย cm แบบ \texttt{float}
มีขั้นตอนการทำงานดังนี้

\begin{lstlisting}[language=C,
    caption={\texttt{hcsr04.c} --- การวัดระยะของ \texttt{HCSR04\_Read()}},
    label={lst:hcsr04_read}]
float HCSR04_Read(void) {
    uint32_t v1, v2, t_start;

    // Trigger pulse
    HAL_GPIO_WritePin(HCSR04_TRIG_PORT, HCSR04_TRIG_PIN, GPIO_PIN_SET);
    __HAL_TIM_SET_COUNTER(hcsr04_htim, 0);
    while (__HAL_TIM_GET_COUNTER(hcsr04_htim) < 10);
    HAL_GPIO_WritePin(HCSR04_TRIG_PORT, HCSR04_TRIG_PIN, GPIO_PIN_RESET);

    // Wait for Echo rising edge
    __HAL_TIM_SET_COUNTER(hcsr04_htim, 0);
    t_start = HAL_GetTick();
    while (HAL_GPIO_ReadPin(HCSR04_ECHO_PORT, HCSR04_ECHO_PIN) == GPIO_PIN_RESET) {
        if ((HAL_GetTick() - t_start) >= HCSR04_ECHO_HIGH_TIMEOUT_MS) return -1.0f;
    }
    v1 = __HAL_TIM_GET_COUNTER(hcsr04_htim);

    // Wait for Echo falling edge
    t_start = HAL_GetTick();
    while (HAL_GPIO_ReadPin(HCSR04_ECHO_PORT, HCSR04_ECHO_PIN) == GPIO_PIN_SET) {
        if ((HAL_GetTick() - t_start) >= HCSR04_ECHO_LOW_TIMEOUT_MS) return -1.0f;
    }
    v2 = __HAL_TIM_GET_COUNTER(hcsr04_htim);

    return (float)(v2 - v1) / 58.0f; // Convert to cm
}
\end{lstlisting}

\subsubsection{Main Loop สำหรับเก็บข้อมูล Raw Data}

Main loop ที่ใช้ในการเก็บข้อมูล Raw Data เพื่อคำนวณ Variance
มีโครงสร้างดังนี้

\begin{lstlisting}[language=C,
    caption={\texttt{main.c} --- Main Loop สำหรับส่ง Raw Distance ออก Serial},
    label={lst:main_raw}]
HAL_TIM_Base_Start(&htim1);
HCSR04_Init(&htim1);

while (1) {
    float raw_cm = HCSR04_Read();

    if (raw_cm > 0.0f) {
        printf("%lu,%.4f\r\n", HAL_GetTick(), raw_cm);
        fflush(stdout);
    }
    HAL_Delay(60);
}
\end{lstlisting}

%% ─────────────────────────────────────────────────────────────────────────────
\subsection{ผลการทดลอง}

\subsubsection{ผลการประเมินค่า Measurement Noise Covariance ($R$)}

จากการเก็บข้อมูล 5 Trials, 500 Samples ต่อ Trial ที่ระยะเป้าหมาย
10, 20, และ 30~cm แล้วคำนวณ Sample Variance ของแต่ละ Trial
และหาค่าเฉลี่ยตามสมการ~\eqref{eq:exp2_mean_variance}
ผลลัพธ์สรุปได้ดังตารางที่~\ref{tab:variance_results}

\begin{table}[htbp]
    \centering
    \caption{ผลการประมวลผล Variance เฉลี่ยจาก 5 Trials,
             500 Samples ต่อ Trial}
    \label{tab:variance_results}
    \begin{tabular}{cccc}
        \hline
        \textbf{ระยะ (cm)}
            & \textbf{ค่าเฉลี่ย $\bar{z}$ (cm)}
            & \textbf{$R = \bar{s}^2$ (cm$^2$)}\\
        \hline
        10 & 9.297  & 0.0050\\
        20 & 18.847 & 0.0337\\
        30 & 28.954 & 0.0734\\
        \hline
    \end{tabular}
\end{table}

จากตารางที่~\ref{tab:variance_results} พบว่าเมื่อระยะห่างของเป้าหมายเพิ่มขึ้น
ความแปรปรวนของสัญญาณก็จะเพิ่มขึ้นตามไปด้วยอย่างชัดเจน

\begin{figure}[htbp]
    \centering
    \includegraphics[width=0.85\textwidth]{exp2_variance_plot.png}
    \caption{กราฟแสดงข้อมูลดิบ (Raw Data) จาก 5 Trials ที่ระยะ 10, 20, และ 30~cm
             แต่ละ Trial บันทึก 500 Samples เพื่อเปรียบเทียบความแตกต่างของ Variance
             ในแต่ละระยะ}
    \label{fig:exp2_variance_plot}
\end{figure}

\subsubsection{การเลือกค่า $R$ สำหรับระบบ Mass-Spring-Damper}

ระบบ Mass-Spring-Damper ในการทดลองถัดไปมีขอบเขตการเคลื่อนที่
สูงสุดประมาณ 20~cm จึงเลือกใช้ $R = 0.0337~\text{cm}^2$
จากระยะ 20~cm เป็นค่าอ้างอิงสำหรับ Kalman Filter ในการทดลองที่ 3 และ 5

% \subsubsection{ขั้นตอนการเลือกค่า \texorpdfstring{$Q$}{Q} สำหรับ Scalar KF}

% การเลือกค่า Process Noise Covariance $Q$ ดำเนินการตามขั้นตอนดังนี้

% \textbf{ขั้นที่ 1 — กำหนดขอบเขตบนและล่างของ $Q$ จากค่า $R$}

% เนื่องจาก $Q$ และ $R$ ควบคุมอัตราส่วนความเชื่อมั่นระหว่าง Model
% และ Sensor ร่วมกัน จึงใช้ $R$ เป็นจุดอ้างอิง:

% \begin{itemize}
%     \item หาก $Q \gg R$: Filter เชื่อ Measurement เกือบทั้งหมด
%           สัญญาณออกจะ Noisy ใกล้เคียง Raw Data
%     \item หาก $Q \ll R$: Filter เชื่อ Model เกือบทั้งหมด
%           สัญญาณออกจะเรียบมากแต่ตอบสนองช้า
%     \item ช่วงที่น่าสนใจจึงอยู่ระหว่าง $Q \approx 10^{-6}$
%           ถึง $Q \approx 10^{-3}$ ซึ่งครอบคลุมทั้งสองด้าน
% \end{itemize}

% \textbf{ขั้นที่ 2 — เลือกค่าเริ่มต้นโดยอ้างอิงจาก Variance ของ Sensor}

% จากการทดลองวัด Variance ของ HC-SR04 ที่ระยะใกล้ (10 cm) ได้
% $s^2 = 0.0050$ cm$^2$ ซึ่งสะท้อนระดับ Noise พื้นฐานที่ต่ำที่สุด
% ในระบบ จึงเลือก $Q = 1 \times 10^{-5}$ เป็นค่า Nominal เนื่องจาก:

% \begin{equation}
%     Q_\text{nominal} = 1 \times 10^{-5}
%     \approx \frac{1}{500} \times R
%     = \frac{1}{500} \times 0.0337
%     \label{eq:Q_nominal_ratio}
% \end{equation}

% ซึ่งทำให้ $Q/R \approx 3 \times 10^{-4}$ กล่าวคือ Filter โน้มเอียง
% ไปทาง Model เป็น Default อันสอดคล้องกับ Static Sensor System
% ที่คาดว่าระยะทางจริงเปลี่ยนแปลงน้อยมากระหว่าง Sampling Period

% \textbf{ขั้นที่ 3 — ตรวจสอบพฤติกรรมจากกราฟ}

% นำค่า $Q_\text{nominal}$ ไปรันกับข้อมูลจริงแล้วสังเกต:

% \begin{itemize}
%     \item สัญญาณออก $\hat{x}_k$ เรียบกว่า Raw Data อย่างเห็นได้ชัด
%           และไม่มี Lag ที่รับรู้ได้ในสภาวะ Static
%     \item เมื่อระยะทางเปลี่ยนแบบ Step Change
%           Filter สามารถติดตามได้ภายในเวลาที่ยอมรับได้
% \end{itemize}

% ผลการทดสอบพบว่า $Q = 1 \times 10^{-5}$ ให้สมดุลที่ดีที่สุด
% ระหว่าง Smoothness และ Responsiveness สำหรับ Application นี้
% สอดคล้องกับค่า Nominal ที่ Al Tahtawi (2018) เลือกใช้เช่นกัน
% ค่า $Q$ ทั้งสามระดับที่ทดสอบในการทดลองที่ 3 จึงเลือกรอบ
% ค่า Nominal นี้ คือ $Q \in \{1\times10^{-3},\,1\times10^{-5},\,1\times10^{-6}\}$

%% ─────────────────────────────────────────────────────────────────────────────
\subsection{สรุปผลการทดลอง}

\begin{itemize}
    \item สามารถพัฒนา Library \texttt{hcsr04.c}/\texttt{hcsr04.h}
          สำหรับอ่านค่าระยะทางจาก HC-SR04 บน STM32G4 ได้สำเร็จ
          โดยใช้ TIM1 เป็น Counter ความละเอียด $1\,\mu$s
          และมีระบบ Timeout เพื่อป้องกันการค้างของโปรแกรม
    \item จากการวัด 5 Trials, 500 Samples ต่อ Trial พบว่า Variance
          ของ HC-SR04 เพิ่มขึ้นตามระยะ: 0.0050~cm$^2$ ที่ 10~cm,
          0.0337~cm$^2$ ที่ 20~cm, และ 0.0734~cm$^2$ ที่ 30~cm
    \item เลือกใช้ $R = 0.0337~\text{cm}^2$ ซึ่งสอดคล้องกับ
          Operating Range ของระบบ Mass-Spring-Damper
          เพื่อนำไปใช้เป็นพารามิเตอร์ใน Kalman Filter ต่อไป
\end{itemize}

%% ─────────────────────────────────────────────────────────────────────────────
\subsection{อภิปรายผล}

\begin{itemize}
    \item \textbf{ปรากฏการณ์การขยายตัวของ Variance ตามระยะทาง:}
          Variance ที่เพิ่มขึ้นจาก 0.0050 ถึง 0.0734~cm$^2$ เป็นผลจาก
          Beam Divergence ของคลื่นอัลตราโซนิค เมื่อเป้าหมายอยู่ไกลขึ้น
          พลังงานคลื่นสะท้อนกลับลดลงและรับสัญญาณรบกวนสิ่งแวดล้อมได้มากขึ้น

    \item \textbf{Systematic Bias ของเซนเซอร์:}
          จากตารางที่~\ref{tab:variance_results} ค่าเฉลี่ยที่วัดได้มีความเบี่ยงเบน
          จากระยะจริง เช่น ที่ระยะ 10~cm วัดได้ 9.297~cm (ต่างจากจริง 0.703~cm)
          ซึ่งเป็น Systematic Error จากการวาง Reference และลักษณะเฉพาะของ Sensor
          ค่า Kalman Filter จะ Converge ไปยังค่าเฉลี่ยนี้ ไม่ใช่ค่าระยะจริง
          จึงต้องคำนึงถึงเมื่อนำไปประยุกต์ใช้

    \item \textbf{ความเหมาะสมในการเลือก $R = 0.0337~\text{cm}^2$:}
          หากใช้ $R$ จากระยะ 30~cm ($= 0.0734~\text{cm}^2$)
          ตัวกรองจะมีความเฉื่อยมากเกินไป ทำให้ตอบสนองช้า
          และหากใช้จาก 10~cm ($= 0.0050~\text{cm}^2$)
          ตัวกรองจะ Noisy มากเกินไป การใช้ค่าจาก 20~cm
          ซึ่งตรงกับ Operating Range จริงจึงเป็น Engineering Tradeoff
          ที่สมเหตุสมผลที่สุด

    \item \textbf{ความน่าเชื่อถือของผล:}
          การวัดซ้ำ 5 Trials ช่วยให้สามารถคำนวณค่า S.D. ของ Variance ข้าม Trials
          เพื่อยืนยัน Reproducibility ของการวัด ค่า S.D. ที่ต่ำแสดงว่า
          Noise Characteristic ของ Sensor มีความสม่ำเสมอและสามารถ
          ใช้เป็นพารามิเตอร์ที่เชื่อถือได้สำหรับ Kalman Filter
\end{itemize}